\documentclass[9pt,twocolumn,twoside,]{pnas-new}

% Use the lineno option to display guide line numbers if required.
% Note that the use of elements such as single-column equations
% may affect the guide line number alignment.


\usepackage[T1]{fontenc}
\usepackage[utf8]{inputenc}

% tightlist command for lists without linebreak
\providecommand{\tightlist}{%
  \setlength{\itemsep}{0pt}\setlength{\parskip}{0pt}}




\templatetype{pnasresearcharticle}  % Choose template

\title{Essai de modélisation de la situation du Chavalier Cuivré
(\emph{Moxostoma hubbsi})}

\author[a,1]{Tristan Plasse}
\author[a,2]{Yohan Wegener}

  \affil[a]{Université de Sherbrooke, Biologie, Sherbrooke, Québec, J1K
2R1}


% Please give the surname of the lead author for the running footer
\leadauthor{}

% Please add here a significance statement to explain the relevance of your work
\significancestatement{}


\authorcontributions{}



\correspondingauthor{\textsuperscript{} }

% Keywords are not mandatory, but authors are strongly encouraged to provide them. If provided, please include two to five keywords, separated by the pipe symbol, e.g:
 \keywords{  Chevalier
Cuivré |  Population |  Menacée |  Ensemencement  } 

\begin{abstract}
Ce travail se voit être une suite logique aux recherches effectuées sur
l'état de la population de Chevalier cuivré (\emph{Moxostoma hubbsi}),
espèce endémique du Québec et considérée menacée. Des estimations
stochastiques basées sur des matrices de leslie à 33 stades distincts
ont été éffectuées par Bannester \& al.~(2025). Nous avons bâtit un
modèle de matrice de levkovitch à 3 stades pour estimer s'il est
possible d'arriver à des résultats aussi concluants avec une méthode
simplifiée. Notre méthode s'est avérée concuante et les résultats
obtenus indiquent les mêmes tendances que ceux obtenus précédement par
Bannester \& al.(2025). L'on peut donc estimer qu'une matrice de
levkovich est suffisante pour modéliser une PVA et qu'une matrice de
leslie n'est pas une nécéssité.
\end{abstract}

\dates{This manuscript was compiled on \today}
\doi{\url{www.pnas.org/cgi/doi/10.1073/pnas.XXXXXXXXXX}}

\begin{document}

% Optional adjustment to line up main text (after abstract) of first page with line numbers, when using both lineno and twocolumn options.
% You should only change this length when you've finalised the article contents.
\verticaladjustment{-2pt}



\maketitle
\thispagestyle{firststyle}
\ifthenelse{\boolean{shortarticle}}{\ifthenelse{\boolean{singlecolumn}}{\abscontentformatted}{\abscontent}}{}

% If your first paragraph (i.e. with the \dropcap) contains a list environment (quote, quotation, theorem, definition, enumerate, itemize...), the line after the list may have some extra indentation. If this is the case, add \parshape=0 to the end of the list environment.

\acknow{}

Le Chevalier cuivré est une espèce de poisson d'eau douce endémique au
Québec. Restreint à certaines zones du fleuve Saint-Laurent et à la
rivière Richelieu, les populations sont précaires. La destruction de son
habitat en est la cause principale.Seulement deux frayères sont connues
à ce jour, et peu de mesures particulières ne sont mises en place pour
les préserver dans le futur. Actuellement, des ensemencements sont
effectuées à chaque année pour tenter de préserver une population
viable. Plus de 14 000 larves et 230 000 alevins sont ainsi enscemencés
dans la rivière Richelieu chaque année pour palier à un taux de
recrutement très faible. Une des principale difficultées concernant la
conservation de cette espèce est le grand manque de connaissances
biologiques sur son histoire de vie. Une étude de XXX a été réalisée
pour estimer la viabilité de la population à long terme de l'espèce
selon 3 scénarios distincs qui prennent en compte la situation actuelle
de l'espèce (source). Cette étude a utilisé un modèle de matrice de
leslie à 33 stades pour effectuer une PVA et pour estimer l'effet de
l'enscemencement à long terme sur la population. Notre objectif pour ce
travail est de répéter leurs analyses tout en construisant un modèle
plus simple de matrice de levkovich à 3 stades seulement.

\section*{Méthodologie}\label{muxe9thodologie}
\addcontentsline{toc}{section}{Méthodologie}

Les données de bases ont directement été tirées de l'article de
Bannester \& al.(2025). Elles se présentes sous la forme d'une matrices
de leslie à 33 satdes où la sous-diagonales représente les survies de
chaque stade et où la fécondité est représentée par la première rangée
(Fig 1). Nous avons choisi de simplifier cette matrice en matrice de
levkovich à 3 stades. Les trois stades intra-annuels de la matrice de
leslie originale, représentés par les termes ``eggs'', ``larvae'' et
``fry'', forment notre premier stade nommé ``Juvéniles''. Les stades 1 à
13, corespondants aux individus adultes non reproducteurs, ont été
agréggés pour former notre deuxième stade nommé ``Jeunes adultes''.
Puis, les stades 14 à 33, correspondants aux individus adultes
reproducteurs, ont été agréggés dans un troisième stade nommé
``Adultes''. La fécondité des individus reproducteurs a été estimée en
faisant la moyenne des fécondités des stades 13 à 14.Comme dans l'étude
de Bannester \& al.(2025), la fécondité du stade reproducteur comporte
une portion de variabilité ajoutée dûe à la stochasticité. Cette
variabilité à été ajouté en assumant que la fécondité annuelle des
individus variait selon une loi normale ayant une moyenne de 0 et un
écart-type de 0,25. Ainsi, la fécondité des individus varie entre aucun
oeuf par année jusqu'au double de la moyenne annuelle. Contrairement à
l'étude de Bannester \& al.(2025),où chaque classe d'âge subissait une
variabilité différente annuellement, tous les individus subissent la
même variabilité au niveau de la fécondité annuelle. Pour la variabilité
concernant les taux de survie, la méthode utilisée est celle de
Bannester \& al.(2025), où l'on assume que ces taux suivent une
distribution beta.

MÉTHODOLOGIE STASES ET SURVIE, MÉTHODOLOGIE STOCHASTICITÉ survie \&
stase.

\subsection*{Résultats}\label{author-affiliations}
\addcontentsline{toc}{subsection}{Résultats}

Include department, institution, and complete address, with the
ZIP/postal code, for each author. Use lower case letters to match
authors with institutions, as shown in the example. Authors with an
ORCID ID may supply this information at submission.

\subsection*{Discussion}\label{submitting-manuscripts}
\addcontentsline{toc}{subsection}{Discussion}

All authors must submit their articles at
\href{http://www.pnascentral.org/cgi-bin/main.plex}{PNAScentral}. If you
are using Overleaf to write your article, you can use the ``Submit to
PNAS'' option in the top bar of the editor window.

\showmatmethods
\showacknow
\pnasbreak



% Bibliography
% \bibliography{pnas-sample}

\end{document}
